%!TEX root = ../../main.tex
\section{교전의 검출과 교전 단위 예측}
\label{sec:rel-encounter}

앞 절의 연구는 경기 전체를 하나의 단위로 보았다. 교전을 단위로 삼으려면 먼저 기록에서 교전을 잘라 내야 한다. 이 절은 교전을 자르는 기존 규칙, 자르는 경계를 정하는 방법, 같은 공개 기록에서 교전 단위의 예측을 먼저 다룬 연구와 그 연구에 남은 한계를 차례로 본다.

\paragraph{교전 검출 규칙.} 교전을 자르는 규칙은 무엇을 전투의 신호로 읽는지에 따라 갈린다. Schubert 등 \cite{schubert2016encounter}은 Dota~2의 게임 기록 전체를 다시 읽어, 상대 유닛이 전투 범위 안에 들어와 두 명 이상이 관여하는 동안을 하나의 encounter로 잘랐다. 그렇게 찾은 encounter의 \killlessSchubert{}\%에는 킬이 없었고, 승자는 결과가 갈린 encounter에서만 예측하였다. Ke 등 \cite{ke2022}은 같은 구조에서 유닛을 잇는 거리를 게임의 공격·치유 사거리로 정하고, 킬이 하나 이상이며 한쪽에 두 명 이상이 있을 것을 요구한 뒤, 찾은 전투를 경기 승자 예측의 입력으로 썼다. Maymin \cite{maymin2021smart}은 리그 오브 레전드에서 챔피언 사이에 피해가 5초 동안 오가지 않으면 교전이 끝난 것으로 보고, 경험치를 나누어 받는 반경 안의 챔피언을 참여자에 더하였다. 이 규칙은 게임 클라이언트에서 얻은 피해 기록이 있어야 쓸 수 있다. 공개 기록 처리기인 OpenDota는 영웅이 죽기 15초 전에 전투 구간을 열고, 죽음 없이 15초가 지나면 닫는다 \cite{opendota2018processteamfights}. Tot 등 \cite{tot2021camera}은 이렇게 표시된 전투의 시각이 일정하지 않다고 지적하였고, Katona 등 \cite{katona2019time}은 죽음을 몇 초 앞서 예측하였다. 이 신호들 가운데 Match-V5의 공개 기록에 시각과 위치가 함께 남는 것은 킬, 곧 챔피언의 죽음뿐이다(제 \ref{sec:data-telemetry} 절). 그래서 본 논문의 규칙은 킬을 기준으로 삼는 쪽에 속한다.

\paragraph{경계의 추정.} 킬을 기준으로 삼으면, 잇따른 킬을 어디까지 한 교전으로 묶을지를 정해야 한다. 같은 꼴의 문제가 다른 분야에 있다. 웹 이용 기록에서 한 사람의 활동을 세션으로 나누는 연구 \cite{halfaker2015session,mehrzadi2012session}와, 지진을 시간과 공간에서 가장 가까운 지진끼리 이어 군집으로 나누는 연구 \cite{zaliapin2013clusters,zaliapin2022perspectives}이다. 두 연구 모두 잇따른 활동 사이의 간격을 모아 분포를 그린다. 같은 묶음 안의 간격은 짧고 묶음 사이의 간격은 길어서 분포에 봉우리가 둘 생기며, 경계는 두 봉우리 사이의 가장 낮은 곳(antimode)에 둔다. 본 논문은 킬 사이의 시간 간격에 같은 방법을 쓴다(제 \ref{sec:eng-temporal} 절).

\paragraph{공개 기록에서의 교전 결과 예측.} 저자와 권준호의 연구 \cite{baek2026killconditioned}(이하 Baek과 Kwon)는 Match-V5의 킬 기록만으로 교전을 정의하고, 그 교전에서 어느 쪽이 이득을 보는지를 예측하였다. 킬은 앞의 킬과 \cogGap{}초 이내이면 한 묶음으로 잇고, 묶음 안에서 가장 먼 두 킬의 거리가 \cogDiam{} 단위를 넘으면 나눈다. 여기서 단위는 게임 지도의 거리 단위이다. 첫 킬 \cogLead{}초 전에 양 팀 각각 \cogGateMin{}명 이상의 살아 있는 챔피언이 첫 킬 위치로부터 \cogGateR{} 단위 안에 있을 때에만 그 묶음을 교전으로 인정한다. 예측에 쓰는 정보는 그 시점까지 \cogWindow{}초 동안의 관측이다. 예측의 정답으로 삼는 양, 곧 라벨은 그 뒤 \cogLabelWindow{}초 안에 오간 교환가치의 부호이다. 교환가치는 킬, 연속 킬, 어시스트, 현상금, 오브젝트, 라인 위치에 손으로 정한 \cogWeights{}개의 가중을 곱해 더한 점수이다. 그 연구는 입력을 표현하는 방식과 학습 알고리즘의 여러 조합을 견주었고, 앞선 패치로 학습한 뒤 학습에 쓰지 않은 나중 패치에서 평가하여 AUC \aucCog{}를 보고하였다. AUC는 이득을 본 교전에 매긴 예측 점수가 그렇지 않은 교전에 매긴 점수보다 높을 확률이며, 0.5가 무작위 수준이다(정의는 제 \ref{sec:rel-scores} 절). 이 값은 가공한 표 형태의 입력에 트리 앙상블인 LightGBM을 쓴 조합에서 나왔고, 시계열과 그래프 구조를 쓰는 심층 모델은 그보다 낮았다. 그 연구는 공개 기록만 쓴다는 제약 아래에서 교전 전 정보가 교전의 결과를 얼마나 알려 주는지를 대규모로 물었고, 교환가치 라벨이 연구자가 정한 조작적 정의임을 스스로 밝혔다. 같은 표 형태의 입력에서 학습 알고리즘을 견준 비교와, 시간순으로 나눈 패치에서의 평가도 포함하였다.

\paragraph{남은 한계.} 이 선행 연구에는 교전의 정의, 결과의 라벨, 예측의 비교의 세 측면에서 추가 검증이 필요한 한계가 남아 있다. 아래에서는 측면마다 한계와 본 논문의 처리를 함께 적는다.

교전의 정의에서는 두 가지가 남는다. 하나는 상수의 근거이다. 킬을 \cogGap{}초로 묶고 \cogDiam{} 단위로 나누는 값은 게임의 성질을 근거로 정해졌고, 자료에서 얻은 근거와 값을 바꾸었을 때의 영향은 보고되지 않았다. 본 논문은 이 두 상수를 자료에서 추정하고, 자료로 추정할 수 없는 상수는 게임 규칙의 값에 맞추어 고정하며, 상수를 바꾸었을 때 교전 집합이 어떻게 달라지는지를 센다(제 \ref{ch:engagement} 장). 다른 하나는 킬이 없는 교전이다. 상대의 기술을 소모시키고 지도 통제권을 가져가면 킬 없이도 교전에서 이길 수 있으므로, 킬로 교전을 찾는 정의는 그런 국면을 분석 대상에서 놓친다. 공개 기록에는 피해 교환이 남지 않으므로 그런 국면 전체의 빈도를 직접 셀 수는 없다. 본 논문은 관측 가능한 대리 조건인 챔피언의 근접을 이용해 킬 없는 근접 국면이 어느 정도 포착되는지를 살펴보고, 그런 사례에서 결과까지 정의할 수 있는지를 검토하여 분석 대상의 범위를 밝힌다(제 \ref{sec:eng-killless} 절).

가장 무거운 한계는 결과의 라벨에 있다. 교환가치는 손으로 정한 가중의 합이고 그 가중의 근거가 약하다. 그래서 무엇을 이긴 교전으로 볼지에 연구자의 판단이 들어가고, 모델이 게임의 흐름이 아니라 라벨을 만든 규칙을 배울 수 있다. 교전의 결과는 전문가도 판정하기 어려운 양이므로, 라벨을 킬 수 우위 같은 다른 결과 정의와 견주고 결과가 라벨의 선택에 얼마나 기대는지를 검토할 필요가 있다. 본 논문은 가중을 손으로 정하지 않는다. 경기 승패로 학습한 뒤 고정한 승률 모형이 매긴 추정 승률이 교전 구간에서 오르는지를 결과로 삼아, 라벨을 경기의 최종 승패와 잇는다(제 \ref{ch:value} 장). 이 결과가 같은 교전의 킬 수 우위와 교환가치 라벨에 얼마나 일치하는지(제 \ref{sec:value-correspondence}, \ref{sec:value-samecase} 절), 승률 모형과 구간을 바꾸면 얼마나 달라지는지(제 \ref{sec:value-horizon} 절)를 함께 보고한다. 이렇게 정한 결과도 승률 모형에 기대어 있으며, 그 한계는 제 \ref{sec:disc-evaluator} 절에서 논의한다.

예측의 비교에서는 두 가지가 남는다. 하나는 학습 알고리즘 비교의 대등성이다. 그 연구의 심층 모델은 규모가 작고 표 형태의 자료를 위한 최근의 심층 모델을 포함하지 않았으며, 가공한 입력은 트리 앙상블 쪽에만 주어졌다. 따라서 그 비교에서 심층 학습의 우열을 읽기는 어렵다. 본 논문은 학습 알고리즘의 우열을 묻지 않는다. 교전 전 상태의 모델을 로지스틱 회귀 하나로 미리 정해 두고(제 \ref{sec:pred-selection} 절), 그 모델이 시작 승률과 경기 시각만 쓰는 기준선보다 덜 틀리는지를 묻는다. 입력의 종류가 더하는 효용은 학습 절차를 하나로 고정한 비교로 따로 재고(제 \ref{sec:pred-infogroups} 절), 같은 입력에서 학습 알고리즘만 바꾸는 비교는 후속 과제로 남긴다(제 \ref{sec:disc-future} 절). 다른 하나는 시간 해상도이다. 관측 구간은 30초인데 상태 기록은 60초마다 남으므로, 구간 안에서 상태가 많아야 한 번 바뀌고 시계열 모델이 쓸 시간 구조가 거의 없다. 본 논문은 이 제약을 자료의 구조로 먼저 밝히고(제 \ref{sec:data-telemetry} 절), 마지막 상태 기록이 얼마나 오래되었는지를 기록하여 검증에서 층화하며, 교전 구간에 새 상태 기록이 들어왔는지에 따라 결과를 나누어 본다(제 \ref{sec:val-frame} 절). 그 시간은 모형 입력에는 넣지 않는다(제 \ref{sec:data-inputs} 절).

\paragraph{측정 단위로서의 교전.} 교전이 무엇인지에 관해 합의된 정의는 없다. 한타는 플레이어의 말이지 측정의 단위가 아니고 \cite{lolwiki2026terminology}, 위의 규칙들도 무엇이 전투를 이루는지를 서로 다르게 본다. 그래서 본 논문은 교전을 단일한 관측 정답이 없는 분석 단위로 두고, 교전을 찾는 규칙을 관측 가능한 기록을 그 단위의 사례로 옮기는 측정 규칙으로 다룬다 \cite{jacobs2021measurement}. 측정 규칙은 가정을 밝혀 적어야 하고, 얼마나 믿을 만하고 타당한지를 검사로 보여야 한다. 본 논문은 Cronbach와 Meehl \cite{cronbach1955construct}의 구성 타당성 개념을 따라 내용 타당성, 수렴 타당성, 판별 타당성, 검사--재검사 신뢰도, 예측 타당성의 증거를 제 \ref{sec:eng-validity} 절에 항목별로 모은다.
